\section*{Exercises}
\begin{exercise}\label{ex:10.1}For $A=\begin{pmatrix}2&1\\0&3\end{pmatrix}$ and $x=(1,-2)$, compute $Ax$. Then explain the answer using the columns of $A$.\end{exercise}
\begin{exercise}\label{ex:10.2}Find the kernel of the linear map $T:\RR^2\to\RR$ given by $T(x,y)=x+y$. Then describe all solutions of $T(x,y)=5$.\end{exercise}
\begin{exercise}\label{ex:10.3}Show that the map $T:\RR^2\to\RR^2$ given by $T(x,y)=(x+y,x-y)$ has inverse $(s,t)\mapsto\bigl((s+t)/2,(s-t)/2\bigr)$. Check both compositions.\end{exercise}
\begin{exercise}\label{ex:10.4}Compute the dot product and the lengths of $u=(1,2)$ and $v=(2,-1)$. Check the Cauchy--Schwarz inequality for this pair.\end{exercise}
\begin{exercise}\label{ex:10.5}Using the Cauchy--Schwarz inequality, prove the triangle inequality $\norm{u+v}_2\le\norm{u}_2+\norm{v}_2$ for all $u,v\in\RR^n$. Begin by comparing the squares of the two sides.\end{exercise}
\begin{exercise}\label{ex:10.6}Let $u_1,\ldots,u_k\in\RR^n$ be nonzero vectors that are pairwise orthogonal, which means that $u_i\cdot u_j=0$ whenever $i\ne j$. Show that $a_1u_1+\cdots+a_ku_k=0$ forces every coefficient $a_j$ to be zero; in other words, the list $u_1,\ldots,u_k$ is linearly independent. Then give an example showing that the word ``nonzero'' cannot be left out.\end{exercise}
